%% ch11 --- Testing: pytest for people who know xUnit, and the trace-assert skeleton
%%
%% Written September 2026 against pytest \pytestver, pytest-asyncio
%% \pytestasyncver, hypothesis \hypothesisver and coverage \coveragever.
%% Every listing is a file under code/ch11/ that runs its own tests when you
%% run it as a script, and CI runs all of them. The three transcripts are
%% generated by code/measure/transcripts.py, two of them from deliberately
%% FAILING runs of the trap_*.py files, which the suite never collects.

\chapter{Testing: pytest for people who know xUnit, and the trace-assert skeleton}
\label{ch:testing}

\noindent
pytest is xUnit with the class removed and the fixture graph made explicit.
That is the whole mapping, and almost everything else follows from it: there
is no base type to inherit, no attribute to attach and no constructor to
overload, so the things those carried \dash{} setup, teardown, shared state,
data-driven cases \dash{} had to go somewhere, and the somewhere is a graph
of named functions that a test asks for by writing their names down.

The parts that transfer are the parts you would expect. The part that does
not is where the incidents come from, and it is not the fixture graph: it is
substituting a dependency, where a Python test can pass for a reason that
has nothing to do with the code under test. That is
\S\ref{sec:testing-patching}, and it is the reason this chapter exists in a
book for people who already know how to test.

\section{What pytest took away}
\label{sec:testing-discovery}
\index{pytest!discovery}

A module whose name begins with \code{test\_}, holding functions whose names
begin with \code{test\_}, is the suite. There is no \code{[Fact]}, no
\code{TestClass} and nothing to register.

\pyregion{code/ch11/test_pricing.py}{plain}{A test is a function in a file, and that is all it is}{lst:ch11-plain}

\noindent
Every listing in this chapter is a file you can run two ways \dash{} under
pytest, or as a script, because of three lines at the bottom of each one:

\pyregion{code/ch11/test_pricing.py}{runner}{The three lines that make a test file runnable on its own}{lst:ch11-runner}

\begin{csbox}
\code{[Fact]} becomes a naming convention; \code{[Theory]} with
\code{[InlineData]} becomes \api{parametrize}
(\S\ref{sec:testing-parametrize}); the constructor and \code{IDisposable}
become a fixture with a \code{yield} in it; \code{IClassFixture<T>} becomes
a scope. The discovery pattern governs \emph{searching}: name a file on the
command line and pytest collects it whatever it is called, which is how the
two deliberately failing files in this chapter are run.
\end{csbox}

\noindent
What replaces \code{Assert.Equal} is stranger, and it is the first thing
worth stopping at. You write a bare \code{assert} \dash{} the language
keyword, the one that vanishes under \code{-O} \dash{} and pytest rewrites
it at import time so that a failure explains itself. Here is what the two
failures in \code{code/ch11/trap\_assert.py} print, generated by running
that file rather than remembered:

\transcript{ch11-assert-rewrite}

\noindent
Nobody wrote a message. pytest recovered the sub-expression that produced
the value, and for the collection it found the index that differs. An
assertion library exists to buy you that, and here the language keyword and
the library are the same thing.

\begin{warning}
Run that same file on a build server and the last line reads differently:
instead of \code{Use -v to get more diff} you get the whole
\pkg{difflib} diff, and long output stops being truncated.
\api{running\_on\_ci} in \pkg{\_pytest.compat} is true when
\code{CI} or \code{BUILD\_NUMBER} is set and non-empty, and the reporting
reads it in more than one place \dash{} the sequence comparison, the
truncation of long output, and the trimming of the short summary line. It
is a deliberate kindness, because on a runner you cannot re-run with
\code{-v}. It also means a report is not reproducible between your machine
and the server unless you say which you want, so the script that generates
every transcript in this book clears both variables.
\end{warning}

\begin{trapbox}
\code{assert} is not the debug-build statement you are thinking of. It is
the assertion library, because pytest rewrites the bytecode of every module
it collects. The consequence is real, though: run the interpreter with
\code{-O} and every \code{assert} in your \emph{application} disappears, so
an \code{assert} is for tests and for invariants you are willing to lose,
never for validating input. Appendix~\ref{app:traps} has this one.
\end{trapbox}

\noindent
Expected exceptions are a context manager rather than an attribute, and the
\code{match} argument is a regular expression against the message:

\pyregion{code/ch11/test_pricing.py}{raises}{Assert.Throws, as a context manager}{lst:ch11-raises}

\section{Fixtures: the constructor, made a graph}
\label{sec:testing-fixtures}
\index{pytest!fixtures}

A fixture is a function with a decorator, and a test asks for one by naming
it as a parameter. Nothing is injected by type, and there is no container:
the parameter name \emph{is} the lookup.

\pyregion{code/ch11/test_pricing.py}{fixture}{A fixture is requested by name, never by type}{lst:ch11-fixture}

\noindent
A fixture may itself request a fixture, which is where the word graph earns
its place \dash{} and where the difference from a constructor stops being
cosmetic. A constructor runs for every test in the class whether the test
wanted it or not; a fixture runs only for the tests that name it, and only
as deep as they name.

\mermaidfig{ch11-fixture-graph}{A test names what it needs, pytest resolves each name to a fixture, and a fixture may name another}{the fixture graph}

\pyregion{code/ch11/test_pricing.py}{graph}{A fixture requesting a fixture}{lst:ch11-graph}

\noindent
Teardown is a \code{yield} rather than a second method. Everything before it
is setup, everything after it runs when the test finishes \dash{} including
when the test fails, which is the guarantee \code{using} gives you and the
reason \code{IAsyncLifetime} has two methods:

\pyregion{code/ch11/test_pricing.py}{teardown}{Setup, yield, teardown, in one function}{lst:ch11-teardown}

\noindent
How long a fixture lives is one argument, and the four values map onto
things you already run:

\begin{center}
\small
\begin{tabularx}{\linewidth}{@{}lX@{}}
\toprule
\textbf{\code{scope=}} & \textbf{What it corresponds to} \\
\midrule
\code{"function"} & The constructor: a new one per test. The default, and
  the one to stay with until something hurts. \\
\code{"class"} & \code{IClassFixture<T>}: shared by the tests in one class,
  in the rare pytest file that has a class at all. \\
\code{"module"} & Shared by one file. \\
\code{"session"} & \code{ICollectionFixture<T>} over everything: built once
  for the whole run. A database container lives here; mutable state
  must not. \\
\bottomrule
\end{tabularx}
\end{center}

\noindent
A fixture that several files need goes in \code{conftest.py} beside them.
That file is not imported by anything and is not named anywhere: pytest
finds it by walking up from each test file, and every fixture in it is
available to everything below. It is the closest thing pytest has to a
composition root, and Chapter~\ref{ch:imports} argued that a composition
root is a function rather than a container.

\begin{trapbox}
There is no test class to put a fixture on, and reaching for one is the
first thing a .NET engineer does. A pytest file with a class in it is
usually a file whose author was still writing xUnit; the class buys grouping
and \code{scope="class"}, and costs an indentation level and a \code{self}
nobody reads. Start without it.
\end{trapbox}

\begin{exercise}{e11_01_fixture}{Remove the duplication with a fixture}
Two tests in the starter build the same basket, line for line. Replace the
duplication with one fixture called \code{basket} and have both tests
request it. The test beside it runs pytest over your file and then looks at
how you wrote it: your tests must pass, the fixture must exist, and the
basket must be built once.
\end{exercise}

\section{One test, four cases}
\label{sec:testing-parametrize}
\index{pytest!parametrize}

\api{parametrize} is \code{[Theory]} and \code{[InlineData]} in one
decorator. The first argument names the parameters, the second is the list
of cases, and pytest generates one test per case with an identifier you can
select on the command line.

\pyregion{code/ch11/test_pricing.py}{parametrize}{Four cases, one test function}{lst:ch11-parametrize}

\begin{csbox}
Each case is a separate test in the report, exactly as an
\code{[InlineData]} row is: four passes rather than one, and a failure names
the case that failed. Where \code{[MemberData]} would take a property,
\api{parametrize} takes any iterable, so a list comprehension or a file read
at import time works without ceremony \dash{} and stacking two
\api{parametrize} decorators on one function gives you the cross product.
\end{csbox}

\begin{exercise}{e11_02_parametrize}{Collapse four tests into one}
The starter has four tests that differ only in their numbers. Make them one
parametrized test that keeps all four cases. The check requires exactly one
test function in the file, the parametrize decorator, and four tests still
passing.
\end{exercise}

\section{Substituting a dependency, and the trap}
\label{sec:testing-patching}
\index{pytest!monkeypatch}

Here is the one thing in this chapter that will cost you an afternoon.

Python has two substitution tools. \api{monkeypatch} is a built-in fixture
that sets an attribute and undoes it at teardown; \pkg{unittest.mock} is the
standard library's Moq, with recording and verification. Both take the thing
to replace as a \emph{string}, and choosing that string wrongly gives you a
test that passes for the wrong reason or fails for no reason you can see.

The module under test does the most ordinary thing in Python:

\begin{python}
# code/ch11/invoice.py
from rates import vat_rate_percent

def gross_bound(net_pence: int) -> int:
    return net_pence + net_pence * vat_rate_percent() // 100
\end{python}

\noindent
So you patch the rate. It is defined in \code{rates}, so that is what you
name. That test is \code{trap\_patch.py}, kept in the repository so that its
failure is a real one:

\transcript{ch11-patch-trap}

\noindent
The patch ran. The rate did not change. Nothing warned.

Chapter~\ref{ch:imports} has the reason, and it is one sentence: a module is
an object that runs once. \code{from rates import vat\_rate\_percent} ran at
import time and bound the function \emph{object} into \code{invoice}'s own
namespace. There are now two names for one function, and rebinding one of
them leaves the other pointing where it always did.

\mermaidfig{ch11-patch-target}{One import, two names for one function object, and only one of them is what the call reads}{where a patch has to land}

\begin{trapbox}
Patch where the name is \emph{looked up}, not where it is defined. Which
name that is was decided by the import: \code{from x import y} binds into
the importing module, so patch \code{importer.y}; \code{import x} followed
by \code{x.y()} looks the attribute up when it calls, so patch \code{x.y}.
The rule is not \enquote{always patch the importer} \dash{} it is
\enquote{patch the name the calling code resolves}. Appendix~\ref{app:traps}
has this one, and it is the entry most worth knowing by heart.
\end{trapbox}

\noindent
Both halves are asserted in \code{code/ch11/test\_patching.py}, so this
book's own build fails the day either stops being true. The right target:

\pyregion{code/ch11/test_patching.py}{right}{Patching the name the call reads}{lst:ch11-right}

\noindent
And the same wrong target is the right one for the other spelling, which is
what makes this a rule about imports rather than a rule about patching:

\pyregion{code/ch11/test_patching.py}{qualified}{A qualified call resolves at call time, so the definition is the target}{lst:ch11-qualified}

\noindent
\pkg{unittest.mock} is the other half. \code{patch} as a context manager is
the scope, and the \code{assert\_called} family is \code{Verify}:

\pyregion{code/ch11/test_patching.py}{mock}{unittest.mock, which is Moq with fewer lambdas}{lst:ch11-mock}

\begin{warning}
A bare \api{Mock} answers to every attribute and accepts every call, so
\code{rate.assert\_called\_onse\_with()} \dash{} note the typo \dash{} is
itself a mock attribute, returns a mock and passes. \api{Mock} has a
\code{spec} argument and \code{patch} has \code{autospec=True} for exactly
this; use one of them wherever the double stands in for something with a
signature worth honouring.
\end{warning}

\begin{exercise}{e11_03_patch_target}{Name the target the call reads}
The starter holds one string: the patch target almost everybody writes
first, which does nothing at all. Change it to the name \code{gross\_bound}
resolves when it calls. The test patches whatever your string names and then
checks that the substitution reached the call.
\end{exercise}

\section{The fixtures you did not write}
\label{sec:testing-builtin}
\index{pytest!built-in fixtures}

pytest ships fixtures for the three things a test most often has to arrange
by hand. \api{tmp\_path} is a \api{Path} to an empty directory, unique to the
test and removed afterwards:

\pyregion{code/ch11/test_builtins.py}{tmp_path}{A scratch directory nobody has to clean up}{lst:ch11-tmppath}

\noindent
\api{caplog} captures log \emph{records} rather than formatted lines, which
is what you want: asserting on the formatted string couples the test to the
formatter, and the formatter is the part most likely to change and least
worth testing.

\pyregion{code/ch11/test_builtins.py}{caplog}{Assert on the record's fields, not on the rendered line}{lst:ch11-caplog}

\noindent
\api{capsys} reads back what was printed, and \code{readouterr()} returns
both streams and clears them.

\pyregion{code/ch11/test_builtins.py}{capsys}{Reading back stdout}{lst:ch11-capsys}

\noindent
Markers are \code{[Trait]}, and selection is \code{-m} for markers and
\code{-k} for a substring of the test's name. A marker has to be registered
in \code{pyproject.toml} or pytest warns \dash{} an unregistered marker is
nearly always a typo, and a typo'd marker is a test that quietly never runs
under the filter meant to select it.

\pyregion{code/ch11/test_builtins.py}{markers}{A marker, and why it has to be declared}{lst:ch11-markers}

\begin{exercise}{e11_04_tmp_path}{Let the test own the directory}
Implement \code{write\_report}, which takes a directory and writes one file
into it. It must not create the directory, clean it up or choose its name:
the test passes \api{tmp\_path}, and a function that decides where its output
goes cannot be tested without touching the machine it runs on.
\end{exercise}

\section{Async tests, property tests, and what coverage measures}
\label{sec:testing-async}
\index{pytest!async tests}

An async test is a test with \code{async} in front of it. There is no
wrapper, no \code{.Result} and no \code{GetAwaiter}, because
\pkg{pytest-asyncio} runs the coroutine on a loop it manages.

\pyregion{code/ch11/test_async_property.py}{async}{An async test, with the loop handled for you}{lst:ch11-async}

\begin{warning}
That works because \code{code/pyproject.toml} sets
\code{asyncio\_mode = "auto"}. In the default strict mode an un-marked async
test is not failed \dash{} it is \emph{skipped}, with a warning. A suite
that reports green having run none of its async tests is the worst outcome
available here, so set the mode once, in the project file, and never per
test. Chapter~\ref{ch:asyncio} is the rest of the story.
\end{warning}

\noindent
hypothesis is FsCheck: state a property, let the library find the case that
breaks it. It also shrinks a failure to its smallest form and remembers it,
so a case found once is replayed on every later run rather than waiting for
the same random draw to come up again.

\pyregion{code/ch11/test_async_property.py}{hypothesis}{A property, and the library invents the inputs}{lst:ch11-hypothesis}

\noindent
Coverage last, and briefly, because the number is worth less than its
reputation. \pkg{coverage} is coverlet: it records which lines executed.
Here is what it says about \code{code/ch11/pricing.py} under a test that
calls every function in it and checks essentially nothing:

\transcript{ch11-coverage}

\noindent
The module is fully covered, in the only sense the tool means. Every line
ran; not one of the assertions in \code{code/ch11/trap\_coverage.py} says
what any of them should have returned, so the module could return any
integer at all and the report would be unchanged. Coverage finds code no
test reaches, which is a real and useful thing to know. It has never been
able to tell you whether the code it reached was checked, and a target
expressed as a percentage buys tests that execute rather than tests that
assert.

\section{trace-assert, stage 01}
\label{sec:testing-traceassert}
\index{trace-assert!stage 01}

\begin{projectbox}
Stage 01: the trace model, the \code{trace} fixture and the first two
assertions, under \code{code/src/trace\_assert/}. Stage 02 records a real
model call into it in Chapter~\ref{ch:aitoolkit};
Chapter~\ref{ch:traceassert} ports the rest of the assertion set, packages
it and compares the three ports.
\end{projectbox}

\noindent
The guiding project is a port of the first layer of agent-eval-bench:
deterministic assertions over what an agent \emph{did}, expressed as pytest
functions. The design is not this book's, and the port follows the original
rather than reinventing it \dash{} three ports of one design is the only
comparison worth making in Chapter~\ref{ch:traceassert}.

Two objects, and the split between them is the whole of stage 01. A
\api{Recorder} is mutable and belongs to the run; a \api{Trace} is immutable
and belongs to the assertions, because an assertion that can edit the
evidence is not an assertion.

\mermaidfig{ch11-trace-assert}{The fixture hands the test a mutable Recorder; the assertions read an immutable snapshot of it}{the trace-assert stage 01 shape}

\noindent
The fixture is four lines in a plugin module, and one line of
\code{conftest.py} registers it. That line becomes an entry point in the
package's own metadata in Chapter~\ref{ch:traceassert}, at which point
installing trace-assert is enough; the staging is deliberate, because an
entry point is packaging.

\pyfile{code/src/trace_assert/plugin.py}{The trace fixture, as a plugin module}{lst:ch11-plugin}

\noindent
Then the two assertions. The schema they come from describes them as a pair:
one asserts that a tool was called, the other that it was not, and a denied
path needs both \dash{} an agent that refuses in prose and calls the tool
anyway passes the first on its own.

\pyfile{code/src/trace_assert/assertions.py}{The first two assertions, and the messages they fail with}{lst:ch11-assertions}

\noindent
Two details in there are pytest rather than trace-assert. \code{\_\_%
tracebackhide\_\_} keeps this file out of the failure report, so the reader
sees the line in \emph{their} test and the message, rather than the line
inside the library that raised. And passing both bounds raises
\api{ValueError} rather than resolving quietly: the original's schema
refuses that combination too, and gives the reason \dash{} one of them would
be ignored, and you would go on believing a bound nothing ever checked.

\begin{exercise}{e11_05_trace}{Assert the refusal and the absence}
A refund agent may pay out up to a limit and must escalate above it.
Implement \code{handle\_refund} so that it records what it did, then read
the test: it asserts the call that should have happened \emph{and} the one
that must not have. Getting the second half right is the exercise.
\end{exercise}

\noindent
That is the suite this book has been running since Chapter~\ref{ch:runtime}:
no network, no database server and no model, which is what makes it able to
run on every push. Evaluating whether an agent's \emph{answer} was any good
is a different problem and belongs to the LangChain volume. What belongs
here is the layer underneath it \dash{} the one that is deterministic, and
therefore the one that can fail a merge.
